% !TEX root = ../main.tex
\section{Elementary Properties of PageRank and Rank Correlation}
\label{ch:appendix-theory}

Chapters~\ref{ch:literature}--\ref{ch:results} rely implicitly on the elementary properties summarized here. The material is standard \cend{15,23,20,21}.

\dissertationSubheading[sec:pr-existence]{Existence and uniqueness of PageRank}

Take $P$ to be a row-stochastic matrix on $n$ nodes, possibly after the dangling-node repair of Equation~\eqref{eq:transition}. For $d\in(0,1)$ the Google matrix
\[
G_{\mathrm{PR}}=d P^{\top}+\frac{1-d}{n}\mathbf{1}\mathbf{1}^{\top}
\]
is positive and, up to the transpose convention of Equation~\eqref{eq:google-matrix}, column-stochastic. By the Perron--Frobenius theorem for positive matrices, $G_{\mathrm{PR}}$ has a single eigenvalue of modulus one, and the corresponding eigenvector $\boldsymbol{\pi}$ is strictly positive; we can scale it so that $\sum_i\pi_i=1$. It follows that, once the damping and the dangling-node rules are fixed, every directed graph has unique PageRank scores \cend{23}.

Uniqueness guarantees a well-defined comparison between EndorseRank and AWP. Once its parameters are fixed, each method assigns exactly one score vector to a given edge set, so any difference between $\mathbf{s}_{\mathrm{ER}}$ and $\mathbf{s}_{\mathrm{AWP}}$ must arise from edge semantics and weights and cannot be traced to multiple solvers.

\dissertationSubheading[sec:pr-continuity]{Continuity in edge weights}

For fixed $d\in(0,1)$, PageRank depends continuously on the entries of $P$. Small perturbations of the allowance and transfer weights thus move the scores only slightly in the $\|\cdot\|_1$ norm, with Lipschitz constants that depend on $d$ and $n$ \cend{23}.

Continuity allows us to run robustness checks that vary the damping or keep only the top tokens. Were alignment overly sensitive to the exact hyperparameters, changes this small would disrupt the family-level $\tau$ altogether. Empirically, EndorseRank's allowance family mean is unchanged in the first three decimal places across $d\in\{0.75,0.85,0.95\}$ (Table~\ref{tab:robustness-damping}), which suggests that the ranking geometry is stable.

\dissertationSubheading[sec:pr-sparsity]{Sparsity and per-iteration cost}

With $m$ nonzero edges, the sparse matrix--vector product $P^{\top}\mathbf{x}$ costs $O(m)$ arithmetic operations, and building the sparse operator from the edge list is also $O(m)$. Power iteration over $I$ iterations is therefore $O(m + mI)$. On the matched cohort the transfer graph has about $9.3$ times as many edges as the endorsement graph, so even at identical $I$ the transfer method would run roughly an order of magnitude longer. Unequal iteration counts widen the gap further, and the measured wall-time ratio (Table~\ref{tab:benchmark-runtime}) is on the same scale as the edge ratio.

\dissertationSubheading[sec:spearman-properties]{Properties of Spearman's \texorpdfstring{$\rho$}{rho}}

Spearman's $\rho$ is Pearson's correlation computed on ranks. It lies in $[-1,1]$ and does not change under strictly increasing transformations of either variable. A value of $\rho=1$ means perfect monotone agreement, and $\rho=-1$ perfect reverse agreement. Tied values receive mid-ranks, and the evaluation pipeline relies on standard tie-aware implementations. Because reputation scores and proxies are continuous and heavy-tailed, rank-based association suits them better than Pearson correlation on the raw values. The AWP baseline paper likewise evaluates its ranking with rank statistics~\cend{3}.

\dissertationSubheading[sec:kendall-properties]{Properties of Kendall's \texorpdfstring{$\tau$}{tau}}

Kendall's $\tau$ is the normalized difference between the numbers of concordant and discordant pairwise orderings. It can also be interpreted probabilistically. For two distinct wallets drawn at random, $\tau$ is the probability that both rankings order them the same way, minus the probability that the rankings order them differently \cend{21}. A $\tau$ in the mid-$0.3$ range, such as EndorseRank attains on the allowance family (Table~\ref{tab:alignment-family-ci}), implies substantial but imperfect agreement between the two rankings.

When $n$ is large, $\tau$ and $\rho$ usually agree in sign and are of similar magnitude, although the two are not identical. We report both, as the AWP baseline does~\cend{3}. Our main summary statistic is the family-mean $\tau$ with a paired-bootstrap interval~\cend{34}.

\dissertationSubheading[sec:construct-validity-note]{Construct validity note}

Cronbach and Meehl~\cend{22} distinguish criterion-related validity from construct validity. Since no definitive gold standard default label exists, this research takes a construct-oriented approach. Each method is tested against validation checks that embody its conceptual definition (e.g.\ receiving allowances for EndorseRank; making inbound transfers for AWP), and also against a holdout period of future new approvals. High within-construct $\tau$ supports the interpretation. A low cross-construct $\tau$, by contrast, does not automatically signal failure; it may indicate discriminant validity.

\dissertationSubheading[sec:numerical-summary-sheet]{Where each primary quantity is reported}

For each primary quantity, the index below names the table where it is established.

\begin{center}
\fitwidth{%
\begin{tabular}{ll}
\toprule
Quantity & Reported in \\
\midrule
Matched-cohort size, $|V|$, $|E|$, runtime, SD, memory, iterations & Table~\ref{tab:benchmark-runtime} \\
In-cohort scaling stages & Table~\ref{tab:benchmark-scaling-incohort} \\
Expanded-pool stages ($N=99{,}080$) & Table~\ref{tab:benchmark-scaling} \\
Family-mean $\tau$ with 95\% CI; inter-method $\tau$ & Table~\ref{tab:alignment-family-ci} \\
Paired $\Delta\tau$ contrasts & Table~\ref{tab:tau-diff} \\
Per-proxy $\rho$, $\tau$, CI & Tables~\ref{tab:alignment-transfer}, \ref{tab:alignment-allowance}, \ref{tab:alignment-sybil-stability} \\
Damping, top-token, sample-definition sweeps & Tables~\ref{tab:robustness-damping}, \ref{tab:robustness-tokens}, \ref{tab:robustness-sample-size} \\
Holdout spender cohort ($n=1{,}335$) & Table~\ref{tab:holdout-spenders} \\
Observation window & 2025-12-01 -- 2026-05-31 (Table~\ref{tab:parameter-summary}) \\
\bottomrule
\end{tabular}
}
\end{center}

Each quantity listed appears both in the evaluation summary and in the tables of Chapter~\ref{ch:results}.
